\documentclass[11pt, a4paper]{article}

\usepackage[utf8]{inputenc}
\usepackage[T1]{fontenc}
\usepackage{amsmath, amssymb, amsfonts}
\usepackage{graphicx}
\usepackage[margin=1in]{geometry}
\usepackage{hyperref}
\usepackage[numbers, sort&compress]{natbib}
\usepackage{booktabs}
\usepackage{siunitx}
\usepackage{caption}
\usepackage{float}
\usepackage{xcolor}

\title{\textbf{Procedural Optimization of Dispersive Readout Pulses with \texttt{cqed\_sim} \texttt{TransmonCavity}}}
\author{Codex \\ cQED Autonomous Research Platform}
\date{\today}

\begin{document}

\maketitle

\begin{abstract}
We study whether a physics-informed, low-dimensional pulse-sequence optimization strategy
can improve dispersive readout in circuit quantum electrodynamics without sacrificing
cavity emptying or repeated-readout quality. The study is implemented on top of
\texttt{cqed\_sim} using the multilevel \texttt{DispersiveTransmonCavityModel} for a
nominal transmon--readout-resonator system with
$\omega_q/2\pi=\SI{6.150}{GHz}$, $\omega_c/2\pi=\SI{8.597}{GHz}$,
$\chi/2\pi=\SI{-2.84}{MHz}$, $\kappa/2\pi=\SI{2.4}{MHz}$, and
$\alpha/2\pi=\SI{-255}{MHz}$. We compare baseline square and smoothed pulses against
three structured alternatives: ring-up/hold/ring-down, free three-segment complex
procedural pulses, analytically nulled tails, and a smooth Fourier basis. A separate
piecewise-constant high-dimensional waveform is used only as an ideal linear reference.

The main result is positive but qualified. In the linear model, the high-dimensional
reference reaches near-perfect detector-limited fidelity by \SI{240}{ns} but leaves
more than five residual photons, so it is not practically acceptable. In the realistic
\texttt{cqed\_sim} replay, structured procedural pulses recover most of that information
while reducing residual photons dramatically. At \SI{240}{ns}, the best free procedural
family reaches $F_\eta=0.956$ compared with $0.871$ for a square pulse, while the
analytic nulling family reaches $F_\eta=0.944$ with only $0.265$ residual photons. By
\SI{496}{ns}, the nulled family reaches $F_\eta=0.9995$ with $3.6\times10^{-3}$ residual
photons. The study therefore supports procedural readout-pulse optimization as a
promising direction for speed and cavity emptying, but not as a complete answer to QND
engineering: within the present two-mode dispersive model, repeated-readout degradation is
dominated by ordinary $T_1$ decay rather than strong-drive state mixing, so a richer model
is still required to resolve the true nonlinear QND limit.
\end{abstract}

\section{Introduction}
Dispersive measurement is the standard readout mechanism for superconducting qubits in
circuit QED \cite{blais2021circuit,koch2007transmon}. In the simplest picture, a driven
readout resonator acquires a qubit-state-dependent response, allowing the qubit state to
be inferred from the transmitted or reflected microwave field
\cite{gambetta2007protocols,boissonneault2009dispersive}. The practical control problem is
not merely to maximize signal separation: useful readout must also leave the resonator
nearly empty, support high repetition rates, and avoid excessive measurement-induced
disturbance \cite{jeffrey2014fast,walter2017rapid,sank2016measurement}.

Existing pulse design is often heuristic. Square or lightly smoothed envelopes are common,
and fast reset strategies such as driven resonator depletion are known to help
\cite{mcclure2016rapid}, but the intermediate regime between fixed heuristics and fully
unrestricted GRAPE-style waveform optimization remains underexplored. That intermediate
regime is precisely what motivates this study: can a small number of interpretable pulse
segments capture most of the useful improvement, especially when the design objective
includes both discrimination and cavity emptying?

The study proceeds in two layers. First, we use a linear dispersive model to define the
objective functions and the bound hierarchy. Second, we replay candidate pulses through the
multilevel \texttt{cqed\_sim} \texttt{DispersiveTransmonCavityModel} with cavity damping
and qubit decoherence. The central question is whether low-dimensional procedural pulses
remain near-optimal after that replay. The remainder of the paper presents the model and
metrics, the numerical results, the validation checks, and the resulting recommendation.

\section{System and Methods}

\subsection{Hamiltonian}
The ideal linear readout bound is described by the state-conditioned resonator equation
\begin{equation}
\dot{\alpha}_{g/e}(t) = -\left( \frac{\kappa}{2} + i \Delta_{g/e} \right)\alpha_{g/e}(t) - i \epsilon(t),
\end{equation}
where $\Delta_g = \omega_c - \omega_d$ and $\Delta_e = \omega_c + \chi - \omega_d$.
The output field is $s_{g/e}(t) = \sqrt{\kappa}\,\alpha_{g/e}(t)$.

The realistic replay uses the two-mode multilevel dispersive Hamiltonian supported by
\texttt{cqed\_sim},
\begin{equation}
\frac{H_0}{\hbar} =
\delta_c a^\dagger a +
\delta_q b^\dagger b +
\frac{\alpha}{2} b^{\dagger 2} b^2 +
\frac{K}{2} a^{\dagger 2} a^2 +
\chi\, a^\dagger a\, b^\dagger b,
\end{equation}
where $a$ is the cavity annihilation operator, $b$ is the transmon lowering operator,
$\delta_c = \omega_c - \omega_c^{\mathrm{frame}}$, and
$\delta_q = \omega_q - \omega_q^{\mathrm{frame}}$. Pulses are applied on the cavity
channel in a frame with $\omega_q^{\mathrm{frame}}=\omega_q$ and
$\omega_c^{\mathrm{frame}}=\omega_d$. Dissipation is included through
\texttt{NoiseSpec(t1, tphi, kappa)}.

\subsection{Simulation Parameters}
\begin{table}[H]
  \centering
  \caption{Nominal system and numerical parameters.}
  \begin{tabular}{@{} l l l @{}}
    \toprule
    Parameter & Value & Unit \\
    \midrule
    Qubit frequency $\omega_q/2\pi$ & 6.150 & GHz \\
    Cavity frequency $\omega_c/2\pi$ & 8.597 & GHz \\
    Anharmonicity $\alpha/2\pi$ & -255 & MHz \\
    Dispersive shift $\chi/2\pi$ & -2.84 & MHz \\
    Cavity linewidth $\kappa/2\pi$ & 2.4 & MHz \\
    Cavity self-Kerr $K/2\pi$ & -28 & kHz \\
    Transmon truncation $n_{\mathrm{tr}}$ & 3 & levels \\
    Cavity truncation $n_{\mathrm{cav}}$ & 14 & Fock states \\
    Detector efficiency $\eta$ & 0.35 & --- \\
    Qubit $T_1$ & 30 & $\mu$s \\
    Qubit $T_\phi$ & 30 & $\mu$s \\
    Time step $\Delta t$ & 4 & ns \\
    Max drive amplitude $|\epsilon_{\max}|/2\pi$ & 4.0 & MHz \\
    Approx. critical photon number $n_{\mathrm{crit}}$ & 20.3 & photons \\
    \bottomrule
  \end{tabular}
\end{table}

\subsection{Computational Approach}
The study-local optimization layer was implemented in
\texttt{studies/procedural\_readout\_pulse\_sequence\_optimization/scripts/readout\_opt/}.
The main reusable components are:
\begin{itemize}
  \item \texttt{pulse\_families.py}: square, smooth square, ring-up/hold/ring-down, free
  three-segment complex pulses, analytically nulled tails, Fourier basis, and an
  8-segment high-dimensional reference.
  \item \texttt{bounds.py}: exact piecewise-constant linear propagation, matched-filter
  SNR, and analytical two-segment cavity nulling.
  \item \texttt{simulate.py}: one- and two-pulse replay through
  \texttt{DispersiveTransmonCavityModel}.
  \item \texttt{optimize.py}: multistart bounded Powell optimization.
\end{itemize}

The primary practical objective was the balanced score
\begin{equation}
J_{\mathrm{bal}} =
F_{\eta}
- 0.08\, n_{\mathrm{res}}
- 0.70\,(1-Q_{\mathrm{QND}})
- 1.20\,L,
\end{equation}
with a peak-photon penalty activated only above $0.45\,n_{\mathrm{crit}}$. Here
$F_\eta$ is the detector-limited matched-filter assignment fidelity,
$n_{\mathrm{res}}$ is the mean residual photon number at the end of the first pulse,
$Q_{\mathrm{QND}}$ is the average state-preservation probability at the start of a second
identical readout, and $L$ is leakage outside the computational subspace. We additionally
used an information-only objective and an emptying-weighted objective to generate a local
tradeoff slice at the representative duration.

Three model gaps had to be documented explicitly. First,
\texttt{ReadoutResonator.response\_trace()} only supports constant drive amplitudes, so
arbitrary-envelope linear bounds were evaluated with study-local exact propagation.
Second, \texttt{cqed\_sim} does not provide a built-in readout-specific optimization
surface, so the objective function is wrapped externally around supported simulator calls.
Third, the available replay path is deterministic Lindblad evolution rather than a full
stochastic measurement unraveling, so repeated-readout quality is quantified by a
two-pulse matched-filter protocol rather than trajectory-conditioned measurement collapse.

\section{Results}

\subsection{Goal 1: Bound hierarchy and the cost of unconstrained optimality}
Figure~\ref{fig:bounds} shows the duration-dependent hierarchy obtained from the
high-dimensional linear reference. The ideal and detector-limited curves reach near-unity
by \SI{240}{ns}; for example, the detector-limited matched-filter fidelity is already
$0.985$ at \SI{240}{ns}. However, this apparent optimum is not practically acceptable: at
the same duration the 8-segment linear reference leaves $n_{\mathrm{res}}=5.57$ photons,
far above the residual-photon budget required for rapid repeated operation.

This is the core structural result of the study. The high-dimensional reference
demonstrates that information extraction itself is not the dominant limitation beyond a few
hundred nanoseconds; instead, the relevant design problem is how to retain most of that
information while enforcing cavity emptying and repeated-use readiness.

\begin{figure}[H]
  \centering
  \includegraphics[width=0.92\linewidth]{../figures/fig1_bound_hierarchy.pdf}
  \caption{Hierarchy of duration-dependent performance curves. The ``$T_1$-limited'' curve
  should be interpreted as a conservative survival-weighted reference, not as a rigorous
  information-theoretic upper bound, because the study does not solve the full decay
  change-point discrimination problem.}
  \label{fig:bounds}
\end{figure}

\subsection{Goal 2: Procedural families versus baseline pulses}
Figure~\ref{fig:frontiers} and Table~\ref{tab:representative} summarize the full-model
family comparison. The baseline square pulse performs poorly at short and intermediate
durations because it spends too much of the allowed pulse window both ringing up and
ringing down passively. At the representative duration of \SI{240}{ns}, the square pulse
reaches only $F_\eta=0.871$ and leaves $1.49$ residual photons.

The structured families do substantially better. At \SI{240}{ns}, the free
three-segment procedural family reaches $F_\eta=0.956$ with $n_{\mathrm{res}}=0.620$,
while the analytic nulling family reaches $F_\eta=0.944$ with only
$n_{\mathrm{res}}=0.265$. The Fourier basis is competitive in fidelity
($F_\eta=0.953$) but not as efficient in emptying. The ring-up/hold/ring-down family is
already strong ($F_\eta=0.947$), which supports the intuition that procedural structure
matters more than unrestricted dimensionality.

\begin{figure}[H]
  \centering
  \includegraphics[width=\linewidth]{../figures/fig2_frontiers.pdf}
  \caption{Full-model frontiers versus duration. The best balanced family is
  \texttt{procedural\_segments} at \SIrange{96}{160}{ns}, \texttt{nulling\_tail} from
  \SIrange{240}{496}{ns}, and a simple ring-up/hold/ring-down pulse at \SI{720}{ns}.}
  \label{fig:frontiers}
\end{figure}

\begin{table}[H]
  \centering
  \caption{Representative \SI{240}{ns} comparison. The piecewise reference is shown only
  for context as an ideal linear benchmark, not as a practical pulse recommendation.}
  \label{tab:representative}
  \begin{tabular}{@{} l c c c c @{}}
    \toprule
    Family & Regime & $F_\eta$ & $n_{\mathrm{res}}$ & $Q_{\mathrm{rep}}$ \\
    \midrule
    Square & Full & 0.871 & 1.485 & 0.853 \\
    Ring-hold & Full & 0.947 & 0.554 & 0.928 \\
    Procedural segments & Full & 0.956 & 0.620 & 0.950 \\
    Nulled tail & Full & 0.944 & 0.265 & 0.936 \\
    Fourier basis & Full & 0.953 & 0.721 & 0.938 \\
    Piecewise reference & Linear & 0.985 & 5.572 & 0.979 \\
    \bottomrule
  \end{tabular}
\end{table}

\subsection{Goal 3: Active cavity emptying}
The analytical nulling result is one of the strongest outcomes of the study. In the linear
model, two free complex tail segments are sufficient to impose
$\alpha_g(T)=\alpha_e(T)=0$ exactly, provided their durations are fixed and the resulting
amplitudes remain within the hardware cap. This gives a concrete answer to one of the main
questions posed at the outset: simultaneous emptying of both state-conditioned cavity
trajectories is not only possible, but algebraically straightforward in the linear model.

In realistic replay the nulling is not exact, but it remains effective. The nulled family
is the best balanced design from \SI{240}{ns} through \SI{496}{ns}. At \SI{336}{ns}, it
reaches $F_\eta=0.990$ with only $4.5\times10^{-2}$ residual photons. At \SI{496}{ns}, it
reaches $F_\eta=0.9995$ with $3.6\times10^{-3}$ residual photons. These results show that
the extra time cost for active emptying is modest once phase variation is allowed.

\begin{figure}[H]
  \centering
  \includegraphics[width=\linewidth]{../figures/fig3_representative_diagnostics.pdf}
  \caption{Representative diagnostics for the free three-segment procedural family:
  output amplitudes for the two basis states, phase-space trajectories
  $(\alpha_g,\alpha_e)$, and the corresponding matched-filter score distributions.}
  \label{fig:diagnostics}
\end{figure}

\subsection{Goal 4: Tradeoffs and robustness}
Figure~\ref{fig:tradeoffs} shows the local tradeoff slice at the representative duration.
Moving from the information-only objective to the emptying-weighted objective pushes the
solutions toward lower residual photon number at a modest fidelity cost. The clearest
beneficiary is the nulled family, which retains high fidelity while moving to substantially
smaller $n_{\mathrm{res}}$.

The robustness study also favors procedural structure. Under the combined perturbation set
($\chi$ scale, $\kappa$ scale, amplitude error, detuning error, and phase error), the
free procedural family retains mean/worst-case fidelities of $0.956/0.945$, compared with
$0.870/0.853$ for the square pulse. The nulled family is slightly less robust than the free
procedural family ($0.944/0.931$) but still decisively better than square.

\begin{figure}[H]
  \centering
  \includegraphics[width=\linewidth]{../figures/fig4_tradeoffs.pdf}
  \caption{Tradeoff slice at the representative duration. Structured families move along a
  favorable fidelity--emptying frontier relative to the baseline square pulse.}
  \label{fig:tradeoffs}
\end{figure}

\begin{figure}[H]
  \centering
  \includegraphics[width=0.82\linewidth]{../figures/fig5_robustness.pdf}
  \caption{Nominal, mean perturbed, and worst perturbed fidelities for representative
  families. The procedural families retain both higher nominal performance and stronger
  robustness than the square pulse.}
  \label{fig:robustness}
\end{figure}

\section{Validation}

\subsection{Sanity Checks}
Three sanity checks were enforced. First, zero drive leaves the cavity empty to numerical
precision. Second, the analytical nulling construction yields
$|\alpha_g(T)|=|\alpha_e(T)|\approx 3.5\times10^{-16}$ in the linear model for a
representative unclipped test case. Third, zero-drive repeated-readout degradation matches
the expected $T_1$ floor rather than a spurious numerical artifact.

\subsection{Convergence Analysis}
The representative full-model procedural-segment pulse was re-evaluated at
$\Delta t=\SI{2}{ns}$ and at cavity truncation $n_{\mathrm{cav}}=18$. The resulting
metric shifts were small:
\begin{align*}
\Delta F_\eta(\Delta t) &= 9.26\times 10^{-4}, \\
\Delta Q_{\mathrm{QND}}(\Delta t) &= 3.29\times 10^{-5}, \\
\Delta n_{\mathrm{res}}(\Delta t) &= 1.07\times 10^{-2}, \\
\Delta F_\eta(n_{\mathrm{cav}}) &= 2.71\times 10^{-5}.
\end{align*}
These shifts are well below the scale of the main family-to-family differences.

\subsection{Literature Comparison}
The study agrees with established dispersive-readout expectations in three ways.
First, the midpoint-drive strategy remains favored, consistent with the standard
continuous-measurement analysis \cite{gambetta2007protocols}. Second, shaped tails are
useful primarily because they accelerate emptying rather than because they change the basic
dispersive measurement geometry, consistent with fast reset/readout literature
\cite{mcclure2016rapid,jeffrey2014fast,walter2017rapid}. Third, the present two-mode
dispersive model does not reproduce strong drive-induced non-QND transitions of the kind
discussed in \cite{sank2016measurement}; that absence is itself a meaningful limitation of
the current modeling layer.

\section{Discussion}
The main scientific conclusion is that the readout optimization problem is well posed in
the present \texttt{cqed\_sim} stack, but only after the objective is defined correctly.
Endpoint separation alone is not sufficient. The practically relevant objective must
combine matched-filter distinguishability with residual-photon and repeated-use penalties.

Within that formulation, low-dimensional procedural control is clearly useful. The free
three-segment family wins the short-time regime, where aggressive loading and extraction
matter most. The nulled family dominates the intermediate regime, where active emptying is
valuable. At long durations the control problem becomes easy enough that a simple
ring-up/hold/ring-down family is sufficient.

The strongest caveat concerns QNDness. The present full-model replay shows essentially no
family-dependent leakage or state-mixing penalty; the reported repeated-readout degradation
is almost entirely the trivial $T_1$ floor over the pulse-plus-wait window. That means the
study \emph{does not} yet establish the ultimate nonlinear QND limit for aggressive
readout pulses. It establishes something narrower but still useful: procedural optimization
is already worthwhile for speed and cavity emptying in the supported two-mode model, while
the unresolved QND question has been localized to missing physics rather than to control
parameterization.

\section{Conclusion}
The answer to the study question is cautiously affirmative. A physics-informed,
procedurally parameterized readout pulse can approach the practically relevant
speed--fidelity frontier much more closely than a square pulse while simultaneously
reducing residual cavity population. In the present implementation:
\begin{itemize}
  \item the optimization problem is well posed in \texttt{cqed\_sim},
  \item matched-filter fidelity plus residual-photon and repeated-readout penalties is the
  correct objective family,
  \item exact simultaneous cavity emptying is achievable in the linear model with two
  complex tail segments,
  \item and low-dimensional procedural families recover most of the practically useful gain
  without the pathological residual occupation of the unrestricted linear reference.
\end{itemize}
The best practical recommendation from this study is therefore not a fully unrestricted
waveform optimizer, but a structured procedural family with complex phase freedom and an
explicit emptying tail. At the same time, answering whether nonlinear backaction truly
kills the advantage requires a richer model than the present two-mode dispersive replay.

\section{Limitations and Future Work}

\subsection{Known Limitations}
The most important limitation is physical rather than numerical: the current
\texttt{DispersiveTransmonCavityModel} replay path is still a two-mode dispersive model.
It captures cavity ring-up, damping, multilevel transmon structure, and ordinary
decoherence, but it does not explicitly generate the stronger measurement-induced state
mixing and dressed-state effects expected deep in the nonlinear regime
\cite{sank2016measurement}. As a result, the family-dependent QND penalty is nearly flat in
the present data.

The plotted ``$T_1$-limited'' curve is a conservative reference, not a rigorous upper
bound. A truly sharp $T_1$-aware limit would require solving a change-point detection
problem for decaying excited-state traces rather than applying a survival-weighted
matched-filter heuristic.

The detector efficiency was treated as a representative fixed parameter
($\eta=0.35$) rather than as a device-calibrated value. This affects the absolute
fidelity scale but not the relative ordering of the pulse families.

\subsection{Suggested Improvements}
\begin{itemize}
  \item \textbf{[P1 | HIGH]} Extend the study to a model with explicit
  readout-induced state-mixing channels or a three-mode readout-aware Hamiltonian. This is
  required to answer the nonlinear-QND question that the current two-mode dispersive model
  leaves unresolved.
  \item \textbf{[P1 | MEDIUM]} Add arbitrary-envelope support to
  \texttt{cqed\_sim.measurement} so that the linear-bound layer and the measurement-chain
  layer use the same public API.
  \item \textbf{[P2 | MEDIUM]} Promote the study-local optimization helpers into a
  reusable readout module once the objective definitions stabilize.
  \item \textbf{[P2 | LOW]} Calibrate the detector efficiency and the cavity self-Kerr from
  device data so the absolute fidelity numbers can be interpreted quantitatively.
  \item \textbf{[P3 | MEDIUM]} Add hardware distortion replay (finite bandwidth, mixer
  skew, and amplitude quantization) to test whether the procedural tails remain robust
  after waveform transport.
\end{itemize}

\subsection{Open Questions}
\begin{itemize}
  \item How much of the procedural-pulse advantage survives once explicit
  measurement-induced transitions are included?
  \item Can the analytically nulled tail remain near-optimal when the tail amplitudes are
  restricted to hardware-realistic slew and bandwidth limits?
  \item Is there a sharper, genuinely rigorous $T_1$-aware readout limit that remains easy
  to evaluate during optimization?
\end{itemize}

\bibliographystyle{unsrtnat}
\bibliography{references}

\end{document}
